\section{Conclusão}
\label{sec:conclusao}

O GEAR organiza uma proposta de governança e gestão de TI para pequenas equipes em três domínios essenciais, um percurso inicial de 30 dias, indicadores e modelos compactos. A construção combina leitura dirigida de referências, síntese e adaptações autorais, com assistência por IA na pesquisa e na produção. O artigo explicita o procedimento de auditoria de conteúdo usado para conferir fontes, preservar limites e corrigir inconsistências.

As contribuições atuais são a especificação documental do framework e o relato rastreável dessa produção assistida. O uso de IA não fornece evidência de eficácia do GEAR, e os registros editoriais não permitem estimar um benefício causal da própria assistência. A ausência de uma revisão sistemática exaustiva, o acesso parcial a algumas fontes e o histórico incompleto de interações delimitam o alcance do relato.

A avaliação do framework permanece futura. O protocolo prevê uma prova de conceito no FrameSim com 25 empresas sintéticas, três configurações pareadas, tarefas comuns e critérios previamente definidos. Nenhuma rodada foi realizada nesta versão e nenhuma tabela contém resultados observados. A prova de conceito poderá revelar cobertura, esforço estimado e falhas nos cenários; evidência de aprendizagem, adesão e efeito organizacional exigirá estudos posteriores em empresas reais.
